% =============================================================================
% Lecture 2: The Foundation Model
% Agentic Machine Learning Systems (Volume III), Chapter 2
% =============================================================================
% Thesis: a model call is a token-bounded proposal across an untrusted boundary,
% so it needs a contract the runtime checks, and its time and price are set
% mostly by output length and by the input the trajectory resends.
% Running case: the repair agent's inspection call, read_file("src/pool.py")
% (Lighthouse 1.1), run for real on a small open model (images/run_model.py).
% Hook and synthesis: War Story 2.1, three serving bugs (Anthropic 2025).
%
% Build: make deck V=3 D=02_foundation_model   (figures: images/make_figs.py)
% Numbers: images/calcs.py prints every derived value; plan in lecture_map.yml.
% =============================================================================

\documentclass[aspectratio=169, 12pt]{beamer}
\usepackage{../../assets/beamerthememlsys}
\mlsyssetup{
  volume={Volume III},
  chapter={Chapter 2},
  chaptertitle={The Foundation Model},
  layer=model,
  logo={../../assets/img/logo-mlsysbook.png},
  instlogo={../../assets/img/logo-harvard.png},
}
\usepackage[T1]{fontenc}
\usepackage[scaled=0.9]{helvet}
\usepackage{courier}
\renewcommand{\familydefault}{\sfdefault}
\usepackage{booktabs,amsmath}
\usetikzlibrary{arrows.meta,positioning}
\graphicspath{{images/}}
\setcounter{mlsystotalsections}{5}

\title{The Foundation Model}
\author{Vijay Janapa Reddi}
\institute{Harvard University}
\date{}

% A large answer choice for predict and concept-question slides
\newcommand{\choice}[2]{\textbf{\textcolor{mlsysaccent}{#1}}\enspace #2}
% A choice that is highlighted from overlay #1 on (the answer)
\newcommand{\answerat}[3]{\alt<#1->{\colorbox{mlsysaccent!15}{\choice{#2}{#3}}}{\choice{#2}{#3}}}
% A goal or takeaway line with its thumbnail
\newcommand{\thumbline}[2]{%
  \begin{minipage}[c]{2.6cm}\centering\includegraphics[width=2.5cm,height=0.95cm,keepaspectratio]{#1}\end{minipage}%
  \hspace{0.3cm}\begin{minipage}[c]{0.74\textwidth}\small #2\end{minipage}}
% One line of a worked example: shown from overlay #1 on
\newcommand{\wstep}[2]{\uncover<#1->{#2}\\[0.2cm]}
% Code and wire text
\newcommand{\cd}[1]{{\ttfamily #1}}
\newcommand{\ind}{\phantom{xx}}

\begin{document}

% =============================================================================
% TITLE AND HOOK
% =============================================================================
\mlsysnote{[0.5 min] [CORE] Read the thesis aloud once. Lecture 1 put a model inside a loop that a runtime closes. Today opens the smallest unit of that loop, one model call, and asks what the runtime must check around it and what it costs. The running case is the repair agent's first inspection, read\_file("src/pool.py"), and we run it for real on a small open model. The notes total about 71 minutes; the core alone runs about 59. FLEX: cut the [OPTIONAL] frames first (the real call, tokenizer, top-p code, kernel bug, truncated call, self-review data, laptop and trace data, open questions); each carries a real artifact whose point the next core frame restates.}
\mlsystitle{The Foundation Model}{A model call is a token-bounded proposal across an untrusted boundary, so the runtime checks it, and output length and resent input set its time and price.}{cover_processor_blueprint.png}

\begin{frame}{Same request, same weights: what would notice?}
  \mlsysactivity{Predict}{3.5 min}
  \centering
  \safeimg[width=0.96\textwidth,height=4.3cm,keepaspectratio]{hook_timeline.pdf}\\[0.05cm]
  {\small Three serving bugs in 2025. Requests and model weights did not change. Green: fix or rollback.}\\[0.12cm]
  {\small\choice{A}{schema check}\quad\choice{B}{HTTP status}\quad\answerat{2}{C}{a test of the output}\quad\choice{D}{temperature 0}}
  \mlsyscredit{Dates and shares redrawn from Anthropic, ``A postmortem of three recent issues,'' September 2025 (anthropic.com/engineering). The Foundation Model, War Story 2.1.}
  \note{[3.5 min] [CORE] Step 1: the dates and shares come from a real public postmortem; the original figure returns at the end of the lecture. Between August and early September 2025, three serving-side bugs degraded a production model's answers. A routing bug sent short requests to servers configured for a much longer context window, first 0.8 percent of one model's requests, then 16 percent at the worst hour; about 30 percent of active coding-agent users met at least one misrouted request. A kernel bug assigned high probability to unlikely tokens, and an approximate top-k selection was miscompiled. Requests and weights did not change. ENGAGE: vote with fingers; which check in an agent's runtime would have noticed? Step 2: C. The corrupted tokens arrived in well-formed responses with clean stop reasons, so the schema check (A) and the transport status (B) passed. Temperature 0 (D) does not help, because the serving path computes the logits. Only a check of the output against the world, a compiler or a test, sees a wrong token. WARN: students who pick A assume malformed text; the failure was fluent text. The synthesis returns to this timeline with the full list of records that help.}
\end{frame}

\begin{frame}{What does a 30-step repair cost to finish?}
  \mlsysactivity{Retrieval}{1.5 min}
  \begin{columns}[c]
    \begin{column}{0.58\textwidth}
      {\Large\bfseries Close your notes.}\\[0.3cm]
      {\normalsize The repair agent takes 30 steps, each correct 99\% of the time. What share of runs finish? What does an accepted repair cost if a \$2.10 attempt is accepted 40\% of the time?}\\[0.3cm]
      \uncover<2->{\small $(1-\epsilon)^{H} = 0.99^{30} = \mathbf{74.0\%}$\qquad $C_{\text{eff}} = \$2.10 / 0.40 = \mathbf{\$5.25}$}
    \end{column}
    \begin{column}{0.38\textwidth}
      \centering\safeimg[width=\textwidth,height=3.4cm,keepaspectratio]{hsa_locator_origin.pdf}\\[0.05cm]
      {\footnotesize One call sits at the origin: one turn, no carried state, zero ambient authority.}
    \end{column}
  \end{columns}
  \mlsyscredit{Lecture 1: horizon compounding (Napkin Math 1.1) and cost per accepted task (Napkin Math 1.2). Locator: The Foundation Model, \S2.1.}
  \note{[1.5 min] [CORE] Retrieval from Lecture 1. Give 45 seconds. Step 2: horizon compounding with per-step error rate $\epsilon$ and horizon $H$ in turns gives 0.99 to the 30th power, 74.0 percent (calcs.py). Cost per accepted task $C_{\text{eff}}$ divides the expected attempt cost by the acceptance probability: \$5.25. Both numbers are properties of the whole trajectory. Point to the locator: H, S, and A are the three exposures from Lecture 1, and a single model call sits at the origin of that space. ENGAGE: if one call is the smallest unit, what can go wrong inside it that the loop must catch? Take two answers and keep them for the synthesis. FLEX: a student who missed Lecture 1 needs only the two formulas on this slide.}
\end{frame}

\begin{frame}{Why does a capable model need a contract per call?}
  \thumbline{thumb_boundary_step3.pdf}{Separate what the model proposes from what the runtime executes, and admit a request by its token count.}\\[0.22cm]
  \thumbline{thumb_probs_step3.pdf}{Show how temperature and top-p reshape candidates, and why only a recorded call replays.}\\[0.22cm]
  \thumbline{thumb_ladder_step4.pdf}{Read the call's status first, and name the one property each check certifies.}\\[0.22cm]
  \thumbline{thumb_deadline_step2.pdf}{Estimate a call's time from its output length, and the input a trajectory resends.}\\[0.12cm]
  {\footnotesize\textcolor{midgray}{Running case: the repair agent's inspection call, read\_file("src/pool.py"), run on a real 0.5B model.}}
  \note{[1 min] [CORE] The driving question from the chapter's Purpose: why does a capable model still need a precise contract around every call? Four goals, one per part of the lecture after the boundary; each thumbnail is the figure that teaches it. The running case is the repair agent from Lecture 1, inspecting src/pool.py before proposing a fix to a connection-pool race. To show real artifacts we ran that exact request on Qwen2.5-Coder-0.5B-Instruct, a small open model, on a laptop (images/run\_model.py). Its outputs are real, and its mistakes are real. WARN: a 0.5B model makes more mistakes than a frontier model; the lecture uses it to show what the runtime receives, not how accurate models are.}
\end{frame}

% =============================================================================
% SECTION: BOUNDARY
% =============================================================================
\mlsysnote{[0.5 min] [CORE] First: who decides what happens to a model's output, and does the request even fit?}
\mlsysdivider{Who acts on a model's output, and does the request fit?}{The Foundation Model, \S2.2 and \S2.3}
\section{Boundary}

\begin{frame}{A model call returns a proposal; the runtime acts on it}
  \centering
  \mlsyssteps[width=0.98\textwidth,height=5.8cm,keepaspectratio]{boundary_step}{3}\\
  {\footnotesize Inference service: the system that hosts the weights and samples tokens.}
  \note{[3 min] [CORE] Step 1: the model and the inference service, the system that hosts the weights and samples tokens, produce only text: read\_file("src/pool.py"). The chapter calls it a proposal (Definition 2.1): it carries zero ambient authority and changes nothing until the runtime validates, authorizes, and executes it. Step 2: the runtime first checks the schema (path is a string, no extra keys), then checks the path against the task's grant, so ../../etc/passwd fails here. Step 3: the tool runs in a sandbox, and the runtime appends the observation, the file excerpt, to the next prompt. The model has no memory between calls; it sees the file only because the runtime resends it. LINK: these are the six phases of a turn from Lecture 1; only the call belongs to the model. ENGAGE: what does a rejected proposal cost? (One model call, counted in $C_{\text{eff}}$; no change to the world.) WARN: logging a tool request as an executed action lets fluent text stand in for evidence. Lecture 7 owns how the runtime mediates every call.}
\end{frame}

\begin{frame}{A real model proposed a well-formed read of the wrong file}
  \begin{columns}[c]
    \begin{column}{0.56\textwidth}
      \begin{mlsyscard}{midgray}\ttfamily\scriptsize
        `{}`{}`json\\
        \{\\
        \ind"name": "read\_file",\\
        \ind"arguments": \{\\
        \ind\ind"path":\\
        \ind\ind\textcolor{RedLine}{\textbf{"/path/to/repair/workspace/file.txt"}}\\
        \ind\}\\
        \}\\
        `{}`{}`<|im\_end|>
      \end{mlsyscard}
      {\footnotesize\textcolor{midgray}{Real output, greedy decoding, ended at end of sequence after 35 tokens. One line break added at the path.}}
    \end{column}
    \begin{column}{0.40\textwidth}
      \small
      \textbf{Asked:} inspect \cd{src/pool.py}\\[0.15cm]
      \textbf{\textcolor{GreenLine}{Valid JSON}} with the right tool name\\[0.1cm]
      \textbf{\textcolor{RedLine}{Wrong path}}, invented by the model\\[0.1cm]
      \textbf{\textcolor{RedLine}{Wrong wrapper}}: a Markdown fence, not the requested \cd{<tool\_call>} tags
    \end{column}
  \end{columns}
  \mlsyscredit{Qwen2.5-Coder-0.5B-Instruct (revision ea3f247), Hugging Face transformers 5.20.1, CPU; request = the chapter's Listing 1. Script: images/run\_model.py.}
  \note{[2 min] [OPTIONAL] This is the real reply to the chapter's Listing 1 request: a system message, the user's request to inspect the locking logic in src/pool.py, and one advertised tool, read\_file with a required string path. The model finished normally after 35 tokens. Its reply parses as JSON and names the right tool. It also names a file that does not exist, and it wraps the call in a Markdown code fence instead of the tool-call tags the chat template asks for. ENGAGE: which runtime stage catches each problem? (A strict tool-call parser catches the wrapper; the schema check passes the path; only authorization and the file system catch the path.) LINK: Lecture 1 called this a fail-plausible fault, output that looks well formed while the intended property fails. WARN: a larger model would more likely name src/pool.py; the runtime must still check, because the same tokens could come from a corrupted serving path, as in the hook.}
\end{frame}

\begin{frame}{The wire request holds 176 tokens; the messages hold 23}
  \begin{columns}[c]
    \begin{column}{0.64\textwidth}
      \begin{mlsyscard}{midgray}\ttfamily\scriptsize
        \textcolor{midgray}{<|im\_start|>system}\\
        \textcolor{mlsysaccent}{\textbf{Propose the next inspection for the}}\\
        \textcolor{mlsysaccent}{\textbf{assigned repair.}}\\
        \textcolor{midgray}{\# Tools}\\
        \textcolor{midgray}{You may call one or more functions [...]}\\
        \textcolor{midgray}{<tools>}\\
        \textcolor{midgray}{\{"type": "function", "function": \{"name":}\\
        \textcolor{midgray}{"read\_file", [...] "required": ["path"] [...]}\\
        \textcolor{midgray}{</tools> [...]<|im\_end|>}\\
        \textcolor{midgray}{<|im\_start|>user}\\
        \textcolor{mlsysaccent}{\textbf{Inspect the connection-pool locking logic}}\\
        \textcolor{mlsysaccent}{\textbf{in src/pool.py.}}\textcolor{midgray}{<|im\_end|>}\\
        \textcolor{midgray}{<|im\_start|>assistant}
      \end{mlsyscard}
    \end{column}
    \begin{column}{0.33\textwidth}
      \centering
      {\Huge\bfseries\textcolor{mlsysaccent}{176}}\\[0.05cm]
      {\small tokens on the wire}\\[0.3cm]
      {\Large\bfseries 23}\\[0.05cm]
      {\small tokens of message text}\\[0.3cm]
      \uncover<2->{\footnotesize Role markers, tool schemas, and format rules: 153 tokens (derived)}
    \end{column}
  \end{columns}
  \mlsyscredit{Real serialized request: Qwen2.5-Coder-0.5B-Instruct chat template with tools=[read\_file] (excerpt; [...] marks elided text; lines wrapped). Counts by the model's own tokenizer. images/run\_model.py.}
  \note{[2 min] [CORE] Step 1: the runtime builds structured messages and tool definitions, but the model receives one serialized sequence. This is the exact wire text for our request, produced by the model's own chat template; gray is framing, purple is the text a person would count. The model's tokenizer counts 176 tokens on the wire against 10 for the system message and 13 for the user message. Step 2: 176 minus 23 is 153 tokens of role markers, tool schema, and format instructions (calcs.py). ENGAGE: the runtime estimates prompt size from the visible messages; by what factor is it wrong here? (7.7 times.) WARN: a remote service may add its own hidden wrapper; then the runtime must use the provider's token-counting endpoint. The admission rule on the next slides counts this whole sequence, never characters or messages. LINK: token-measured admission is this lecture's; Lecture 4 decides what goes into the working set.}
\end{frame}

\begin{frame}{Four spaces of indentation change the token sequence}
  \centering
  \safeimg[width=0.94\textwidth,height=4.2cm,keepaspectratio]{tokens.pdf}\\[0.15cm]
  {\small \cd{src/pool.py} is one argument at the tool and 4 tokens at the model: \cd{src}, \cd{/p}, \cd{ool}, \cd{.py}.}
  \mlsyscredit{Real tokenizer output: tiktoken 0.14.0, cl100k\_base byte-pair encoding (BPE). The Foundation Model, \S2.3.1 and figure 2.}
  \note{[1.5 min] [OPTIONAL] Real tokenizer output from tiktoken's cl100k\_base vocabulary, the example the chapter uses (\S2.3.1). Byte-pair encoding, BPE, learns frequent fragments from a training corpus, so token boundaries rarely match words or identifiers. The flush line def get\_user\_id(): is 5 tokens, starting with def at id 755. Indented by four spaces it is 6 tokens: three spaces (262), then " def" with the fourth space fused to it (711), then the same four tokens. A compiler's lexer sees the same keyword at both indents; the model sees a different opening sequence. The path src/pool.py is 4 tokens. ENGAGE: why can a runtime not use a fixed characters-per-token ratio? (Indentation, escaping, and rare identifiers change the count unpredictably.) WARN: counts differ across model families; this vocabulary is not the 0.5B model's. FLEX: skip if short on time; the previous slide carries the rule.}
\end{frame}

\begin{frame}{Admission counts the prompt and reserves the reply}
  \centering
  {\large Token-measured admission\qquad $S \;\le\; S_{\max} - K_{\max}$}\\[0.05cm]
  {\footnotesize $S$: serialized prompt tokens\quad $S_{\max}$: context window\quad $K_{\max}$: reserved output tokens}\\[0.1cm]
  \mlsyssteps[width=0.96\textwidth,height=4.8cm,keepaspectratio]{admit_step}{3}
  \mlsyscredit{The Foundation Model, equation 1, \S2.3.2 and \S2.3.3 (scenario values).}
  \note{[2.5 min] [CORE] Equation 1: a request is admitted only if its serialized length $S$ fits in the context window $S_{\max}$ with $K_{\max}$ tokens reserved for the reply. Step 1: the chapter's scenario, a window of 8,192 tokens and a reserve of 1,024, so the input ceiling is 7,168; a 6,000-token prompt leaves 1,168 to spare. Step 2: attach a screenshot of the failing dashboard. The vision front end reports 2,048 more context positions, so demand is 9,072, 880 over the window, and admission fails before the call is sent. Step 3: the next slide's question. WARN: reasoning models spend hidden reasoning tokens before the visible reply, and those count against $K_{\max}$; a three-line answer can exhaust a generous reserve. ENGAGE: what can the runtime do instead of failing? (Compress history, take a narrower file slice, downscale the image, or route to a larger window.) LINK: Lecture 5 prices what those tokens hold in memory.}
\end{frame}

\begin{frame}{Does halving the reply reserve admit the image?}
  \mlsysactivity{Estimate}{2.5 min}
  \begin{columns}[T]
    \begin{column}{0.40\textwidth}
      \begin{mlsyscard}{mlsysaccent}
        \small\textbf{Given (Checkpoint 2.1)}\\[0.05cm]
        Window: 8,192\\
        Prompt: 6,000\\
        Image: 2,048\\
        Reserve: 1,024 $\rightarrow$ 512
      \end{mlsyscard}
    \end{column}
    \begin{column}{0.56\textwidth}
      \small
      \wstep{2}{$6{,}000 + 2{,}048 + 512 = 8{,}560$}
      \wstep{2}{$8{,}560 - 8{,}192 = \mathbf{\textcolor{RedLine}{368}}$ tokens over (derived)}
      \uncover<3->{The request still fails, and a 512-token reserve can cut a tool call mid-argument. The reply then ends \cd{TRUNCATED}.}
    \end{column}
  \end{columns}
  \mlsyscredit{The Foundation Model, Checkpoint 2.1. Arithmetic: images/calcs.py.}
  \note{[2.5 min] [CORE] ENGAGE: one minute alone, then compare with a neighbor. A teammate proposes halving the output reserve instead of shrinking the image. Does the request fit? Step 2: 6,000 plus 2,048 plus 512 is 8,560, still 368 over. Step 3: and even if it fit, a smaller reserve risks a reply that stops mid-argument; the next segment shows what that reply looks like. Expected answers: most students compute 368 correctly; the common error is to subtract only the image. WARN: shrinking the reserve to force admission trades a visible failure now for an invisible one later. The better moves are a narrower file slice, a downscaled image (at the cost of blurred error codes), or a larger window.}
\end{frame}

% =============================================================================
% SECTION: GENERATION
% =============================================================================
\mlsysnote{[0.5 min] [CORE] Land the first segment aloud: the model proposes, the runtime decides, and a request is admitted by its full token count with the reply reserved. Now: how does the service turn a fixed request into tokens, and why do repeated calls differ?}
\mlsysdivider{Why do repeated calls differ, and what replays exactly?}{The Foundation Model, \S2.4}
\section{Generation}

\begin{frame}{Temperature sharpens the guess the model already prefers}
  \centering
  \mlsyssteps[width=0.96\textwidth,height=5.0cm,keepaspectratio]{probs_step}{3}\\[0.02cm]
  {\footnotesize Next token after \cd{"path": "}. Sampling shapes candidates; only a check decides correctness.}
  \mlsyscredit{Real next-token probabilities: Qwen2.5-Coder-0.5B-Instruct on the chapter's Listing 1 request (images/run\_model.py; top 8 of 151,936 tokens).}
  \note{[3 min] [CORE] Real data from our run. At each step the model scores every token in its 151,936-token vocabulary; the scores are logits, and softmax with temperature $T$ (equation 3) turns them into probabilities. Step 1: at the first token of the path, at $T = 1$, the model puts 34.6 percent on "path" (the start of /path/to/...), 16.3 percent on "C", and only 11.3 percent on "src", the file the prompt names. Greedy decoding takes "path". Step 2: halving the temperature to 0.5 raises "path" to 69.2 percent and cuts "src" to 7.4 percent. Lower temperature concentrates mass on the token the model already prefers; it does not move mass toward the correct one. Step 3: top-p 0.90 keeps the shortest set of tokens whose mass reaches 0.90; here that is 32 tokens, because this distribution is flat. ENGAGE: which setting is "more accurate"? (Neither; only a pass rate on the task can say.) WARN: high probability is not confidence in correctness; training rewards predicting text, not finishing tasks.}
\end{frame}

\begin{frame}{How much mass does halving temperature move?}
  \mlsysactivity{Predict}{2.5 min}
  \begin{columns}[c]
    \begin{column}{0.56\textwidth}
      \centering
      \mlsyssteps[width=\textwidth,height=4.85cm,keepaspectratio]{nm21_step}{2}\\
      {\footnotesize Bars: tokens $v_1$ to $v_5$, labeled by logit.}
    \end{column}
    \begin{column}{0.40\textwidth}
      \small Logits $[12, 10.5, 9, 6, 3]$; $p(v_1) = 78.4\%$ at $T = 1$.\\[0.15cm]
      \textbf{At $T = 0.5$, what does $v_1$ get?}\\[0.1cm]
      \choice{A}{80\%}\quad\answerat{2}{B}{95\%}\quad\choice{C}{100\%}\\[0.2cm]
      \uncover<2->{\footnotesize Tail: 21.6\% $\rightarrow$ 5.0\%, a 4.3$\times$ cut. Top-p 0.90 at $T = 1$ keeps $v_1, v_2$ (0.9589) and renormalizes to 81.8\% and 18.2\%.}
    \end{column}
  \end{columns}
  \mlsyscredit{The Foundation Model, Napkin Math 2.1 (hypothetical logits). Arithmetic: images/calcs.py.}
  \note{[2.5 min] [CORE] The chapter's worked example, with logits held fixed so only the decoding setting changes. ENGAGE: vote. Halving $T$ doubles every exponent. Step 2: B, 95.0 percent. The non-top tail falls from 21.6 to 5.0 percent, a 4.3 times cut, and no token reaches zero. Students who choose C confuse low temperature with greedy decoding; only the limit $T \rightarrow 0$, or an explicit argmax, removes the tail. Students who choose A underestimate the exponential. Top-p 0.90 at $T = 1$: $v_1$ alone has 0.7840, short of 0.90; adding $v_2$ reaches 0.9589, so the nucleus is two tokens, renormalized to 0.8176 and 0.1824. Compare the real call on the previous slide, where top-p 0.90 kept 32 tokens: top-p adapts to the spread. LINK: the tail is the diversity Lecture 3 spends when it samples many candidates.}
\end{frame}

\begin{frame}{Top-p sampling is twelve lines that never see the task}
  \begin{columns}[c]
    \begin{column}{0.76\textwidth}
      \begin{mlsyscard}{midgray}\ttfamily\scriptsize
        def \_\_call\_\_(self, input\_ids, scores):\\
        \ind sorted\_logits, sorted\_indices = torch.sort(\\
        \ind\ind scores, descending=False)\\
        \ind cumulative\_probs = sorted\_logits.softmax(\\
        \ind\ind dim=-1).cumsum(dim=-1)\\
        \ind sorted\_indices\_to\_remove = (\\
        \ind\ind cumulative\_probs <= (1 - self.top\_p))\\
        \ind sorted\_indices\_to\_remove[\\
        \ind\ind ..., -self.min\_tokens\_to\_keep :] = 0\\
        \ind indices\_to\_remove = sorted\_indices\_to\_remove\\
        \ind\ind .scatter(1, sorted\_indices, sorted\_indices\_to\_remove)\\
        \ind scores\_processed = scores.\textbf{\textcolor{RedLine}{masked\_fill}}(\\
        \ind\ind indices\_to\_remove, self.filter\_value)\\
        \ind return scores\_processed
      \end{mlsyscard}
    \end{column}
    \begin{column}{0.21\textwidth}
      \small Inputs: scores and a threshold.\\[0.15cm]
      No file system, no grant, no test.\\[0.15cm]
      \textcolor{midgray}{\footnotesize Type hints and two comments dropped; lines wrapped to fit.}
    \end{column}
  \end{columns}
  \mlsyscredit{Hugging Face transformers 5.20.1, src/transformers/generation/logits\_process.py, TopPLogitsWarper.\_\_call\_\_ (lines 543 to 555).}
  \note{[1 min] [OPTIONAL] The real top-p filter from a widely used framework. It sorts the scores, accumulates probability, marks the low-probability tail, and fills those scores with negative infinity so the sampler never draws them. Its only inputs are the scores and a threshold. Nothing here can know that src/pool.py exists or that /path/to/... does not. ENGAGE: where in the runtime does that knowledge enter? (Authorization and the tool observation, after the call.) WARN: the code sorts ascending and removes the tail below $1 - p$; it is the same rule as the slide before, written from the other end. FLEX: skip if short on time.}
\end{frame}

\begin{frame}{Sharing a batch changed 98.8\% of one request's logits}
  \begin{columns}[c]
    \begin{column}{0.50\textwidth}
      \small Same request, run alone and then batched with a 471-token request (295 padding tokens added):\\[0.25cm]
      \begin{tabular}{@{}lrr@{}}\toprule
        & \textbf{32-bit} & \textbf{16-bit (bf16)} \\\midrule
        Logits changed & 150,078 & 140,555 \\
        Largest change & $3.9 \times 10^{-5}$ & 0.25 \\
        Top token & same & same \\\bottomrule
      \end{tabular}\\[0.2cm]
      {\footnotesize Of 151,936 logits. bf16: a 16-bit float with 8 exponent bits.}
    \end{column}
    \begin{column}{0.46\textwidth}
      \centering
      \uncover<2->{\begin{mlsyscard}{mlsysaccent}\small
        A near tie at one position flips the argmax, and every later token follows the new prefix.\\[0.1cm]
        \textbf{Log} model revision, serving settings, and sampling settings; \textbf{replay} the recorded response.
      \end{mlsyscard}}
    \end{column}
  \end{columns}
  \mlsyscredit{Measured: Qwen2.5-Coder-0.5B-Instruct, PyTorch 2.14.1 CPU, last-position logits alone versus left-padded batch of two (images/run\_model.py). The Foundation Model, \S2.4.3.}
  \note{[2 min] [CORE] A real measurement. We computed the next-token logits for our request alone, then again inside a batch of two with a longer request, which pads ours with 295 tokens. In 32-bit floats, 150,078 of 151,936 logits changed (98.8 percent), by at most $3.9 \times 10^{-5}$; in bf16, 140,555 changed, by up to 0.25. The top token stayed the same here. Step 2: floating-point addition is not associative, so a different batch shape sums the same numbers in a different order (\S2.4.3). When two tokens nearly tie, that is enough to flip a greedy choice, and every later token then differs. ENGAGE: does fixing the random seed prevent this? (No; the seed controls the sampler, not the arithmetic.) WARN: in production the batch depends on other users' traffic, which the caller cannot see. LINK: Lecture 11 builds record and replay into the runtime; today the rule is only that a recorded response replays and a served call does not.}
\end{frame}

\begin{frame}{A miscompiled top-k kernel dropped the largest logit}
  \begin{columns}[c]
    \begin{column}{0.62\textwidth}
      \begin{mlsyscard}{midgray}\ttfamily\scriptsize
        k = 256\\
        N = 12000\\
        for i in range(N):\\
        \ind arr = jnp.zeros((N,))\\
        \ind arr = arr.at[i].set(1.0)\\
        \ind approx\_values, approx\_indices = \\
        \ind\ind lax.\textbf{\textcolor{RedLine}{approx\_max\_k}}(arr, k=k)\\
        \ind exact\_values, exact\_indices = \\
        \ind\ind lax.top\_k(arr, k=k)\\
        \ind if approx\_values[0] != exact\_values[0]:\\
        \ind\ind print(f"Diff at \{N=\} \{i=\} [...]")
      \end{mlsyscard}
    \end{column}
    \begin{column}{0.35\textwidth}
      \small An array of 12,000 zeros with a single one. On the accelerator, the approximate top-k missed the maximum for every index $\ge$ 10,240.\\[0.2cm]
      Fix: exact top-k in 32-bit floats, with a small throughput cost.
    \end{column}
  \end{columns}
  \mlsyscredit{Reproduction code as published in Anthropic, ``A postmortem of three recent issues,'' 2025 (imports and two comments omitted; lines wrapped). The Foundation Model, War Story 2.1.}
  \note{[1.5 min] [OPTIONAL] The third bug from the hook, in the code the engineers used to reproduce it, published in the postmortem. The test builds an array of 12,000 zeros with a single one and asks both the approximate and the exact top-k for the largest value. On the accelerator the approximate routine failed for every position at or after 10,240. In a sampler, that means the most likely token can silently drop out of the candidate set. The fix replaced it with an exact top-k in 32-bit precision, trading a little throughput for correct selection. ENGAGE: would a caller's logs show this? (Only if they record the serving configuration and test outputs against the world; the request and the response format looked normal.) WARN: this is a serving-path bug, not a model bug; retraining would not fix it. FLEX: cut if short on time.}
\end{frame}

\begin{frame}{True or false: greedy decoding replays a call?}
  \mlsysactivity{Concept question}{2 min}
  \centering
  {\normalsize ``We set temperature to 0, so a failed agent run can be reproduced by sending the same request again.''}\\[0.4cm]
  {\Large\choice{T}{True}\qquad\answerat{2}{F}{False}}\\[0.4cm]
  \uncover<2->{\begin{mlsyscard}{mlsysaccent}\small
    Greedy fixes the choice for given logits, but the serving path computes the logits. Batch shape (98.8\% of logits shifted) and serving changes (16\% of requests at peak) both move them.
  \end{mlsyscard}}
  \mlsyscredit{The Foundation Model, \S2.10 (fallacy: greedy decoding makes a served call reproducible).}
  \note{[2 min] [CORE] The chapter's fallacy as a concept question. Vote, discuss with a neighbor for 30 seconds, vote again. Step 2: false. Greedy selection is deterministic for a fixed logit vector and tie rule, but two causes change the vector. On a healthy serving path, batch-dependent reduction order shifts logits (our measurement, 98.8 percent of logits moved). A change in the serving path itself, as in the hook, changed outputs for identical requests (16 percent of one model's requests at the worst hour). ENGAGE: what should the runtime log with every call? (The serialized prompt, model and tokenizer identifiers, sampling settings, token limits, routing or seed metadata, the raw response, its stop reason, and token usage.) Then replay the recorded response to test downstream logic, and resend the live request only to measure variation. WARN: students who vote true are thinking of a single machine running one request.}
\end{frame}

% =============================================================================
% SECTION: CONTRACT
% =============================================================================
\mlsysnote{[0.5 min] [CORE] Land the second segment: sampling reshapes candidates without testing them, and only a recorded call replays exactly. Now: how does the runtime know a complete candidate arrived at all?}
\mlsysdivider{How does the runtime know a whole candidate arrived?}{The Foundation Model, \S2.5}
\section{Contract}

\begin{frame}{A truncated command can be valid and delete more}
  \centering
  \safeimg[width=0.86\textwidth,height=4.6cm,keepaspectratio]{rmrf.pdf}\\[0.1cm]
  {\small Both strings are valid shell commands. Only the call's status tells them apart.}
  \mlsyscredit{The Foundation Model, \S2.5 and \S2.5.2.}
  \note{[1.5 min] [CORE] The chapter's example. The model starts a cleanup command for /tmp/scratch/cache\_dir and hits its token limit right after /tmp/scratch. The cut-off prefix is a valid shell command, and it deletes the whole parent directory: build outputs, results, everything beside the cache. A runtime that reads only the text cannot tell an intended short command from a cut one. ENGAGE: which field must arrive with the text to tell them apart? (How the call ended: completed, or stopped at the limit.) WARN: students propose checking whether the text parses; both strings parse. This is why the invocation contract pairs every payload with a status.}
\end{frame}

\begin{frame}{The runtime reads the status before the payload}
  \begin{columns}[c]
    \begin{column}{0.58\textwidth}
      \footnotesize
      \begin{tabular}{@{}P{3.4cm}P{4.4cm}@{}}\toprule
        \textbf{Status} & \textbf{Host action} \\\midrule
        \textcolor{GreenLine}{\textbf{COMPLETED}} & admit to candidate checks, not to dispatch \\
        \textcolor{RedLine}{\textbf{TRUNCATED}} & withhold; record the limit \\
        REFUSED & apply the refusal policy \\
        TRANSPORT\_FAILURE & remote outcome unknown \\
        REJECTED, CANCELED & stop locally; track the remote \\
        FAULTED & route to failure handling \\\bottomrule
      \end{tabular}
    \end{column}
    \begin{column}{0.39\textwidth}
      \uncover<2->{\begin{mlsyscard}{midgray}\ttfamily\scriptsize
        def completed\_candidate(env):\\
        \ind if env["status"] !=\\
        \ind\ind\ind "COMPLETED":\\
        \ind\ind raise ValueError(\\
        \ind\ind\ind "not a completed\\
        \ind\ind\ind candidate")\\
        \ind return env["payload"]
      \end{mlsyscard}
      {\footnotesize The invocation contract pairs each request with this status. Only COMPLETED enters the candidate path.}}
    \end{column}
  \end{columns}
  \mlsyscredit{The Foundation Model, Definition 2.2, table 2, and Listing 3 (variable renamed and message shortened to fit).}
  \note{[3 min] [CORE] The invocation contract (Definition 2.2) is a pair of records: the request (prompt, model id, output limit, deadline, stop sequences, sampling settings) and a normalized envelope (status, payload, metadata with stop reason and token usage). Step 1: the runtime adapter maps every provider's signals onto seven statuses. COMPLETED means generation ended normally within limits, and nothing more: syntax, permission, and truth are still unchecked. TRANSPORT\_FAILURE records an unknown remote outcome, because a local timeout cannot tell a lost request from a lost reply (Birrell and Nelson's remote procedure call problem); the call may have run and been billed. Step 2: the quarantining rule as code. Anything not COMPLETED goes to a diagnostic log and never to dispatch. Streaming lets a monitor reject a bad prefix early: rejected at step 96 of 4,096, an honored cancel skips 4,000 steps, 84.3 s at the reference decode floor, but only if the server honors the cancel. WARN: repairing a fragment by closing its braces invents a call the model never made. LINK: Lecture 12 decides what to do next for each status.}
\end{frame}

\begin{frame}{A real call cut at 24 tokens still looks like a tool call}
  \begin{columns}[T]
    \begin{column}{0.48\textwidth}
      {\small\textbf{max\_new\_tokens = 128}}\\[0.05cm]
      \begin{mlsyscard}{GreenLine}\ttfamily\scriptsize
        \{\\
        \ind"name": "read\_file",\\
        \ind"arguments": \{\\
        \ind\ind"path":\\
        "/path/to/repair/workspace/file.txt"\\
        \ind\}\\
        \}
      \end{mlsyscard}
      {\footnotesize status COMPLETED (end of sequence), 35 tokens}
    \end{column}
    \begin{column}{0.48\textwidth}
      {\small\textbf{max\_new\_tokens = 24}}\\[0.05cm]
      \begin{mlsyscard}{RedLine}\ttfamily\scriptsize
        \{\\
        \ind"name": "read\_file",\\
        \ind"arguments": \{\\
        \ind\ind"path":\\
        \ind\ind"/path/to
      \end{mlsyscard}
      {\footnotesize status TRUNCATED (length), 24 tokens}\\[0.1cm]
      \uncover<2->{\footnotesize Closing the braces would invent \cd{read\_file("/path/to")}.}
    \end{column}
  \end{columns}
  \mlsyscredit{Real outputs, same request, greedy, two output limits (fence lines omitted; break added at the path). Qwen2.5-Coder-0.5B-Instruct; images/run\_model.py.}
  \note{[1.5 min] [OPTIONAL] Two real calls with the same request, differing only in the output limit. With 128 tokens the model finishes after 35 and the engine reports end of sequence: status COMPLETED. With 24 tokens it stops mid-path and the engine reports the length limit: status TRUNCATED. The first three lines are identical, and the fragment is a legal JSON prefix. Step 2: a lenient "repair parser" that closes the string and braces would produce read\_file("/path/to"), a request no one made, for a directory instead of a file. ENGAGE: which record, kept with the payload, prevents that? (The stop reason, mapped to TRUNCATED.) WARN: retrying the same request at the same limit truncates again; the fix is a larger allowance or a more compact output format.}
\end{frame}

\begin{frame}{Which status, and which next action?}
  \mlsysactivity{Diagnose}{3 min}
  \footnotesize
  \begin{tabular}{@{}P{5.0cm}P{3.5cm}P{4.8cm}@{}}\toprule
    \textbf{Symptom} & \textbf{Status} & \textbf{Next action} \\\midrule
    1. The 15 s deadline passes with no reply & \uncover<2->{TRANSPORT\_FAILURE} & \uncover<2->{record unknown; check before resending} \\
    2. The reply stops at max\_output\_tokens & \uncover<2->{TRUNCATED} & \uncover<2->{larger allowance or shorter format} \\
    3. A safety refusal, no tool call & \uncover<2->{REFUSED} & \uncover<2->{refusal policy; do not retry blindly} \\
    4. Valid JSON naming a missing file & \uncover<2->{COMPLETED} & \uncover<2->{authorization fails; read the directory} \\\bottomrule
  \end{tabular}\\[0.1cm]
  \uncover<3->{\footnotesize A blind retry repeats the cause in rows 2 to 4. In row 1 the first call may still be running and billed.}
  \mlsyscredit{The Foundation Model, table 2 and \S2.10 (pitfall: collapsing every unsuccessful outcome into the same retry).}
  \note{[3 min] [CORE] ENGAGE: pairs, two minutes. For each symptom name the status the adapter records and the next action. Step 2: (1) a local deadline expiry is TRANSPORT\_FAILURE with an unknown remote outcome; the server may have finished and billed the call. (2) TRUNCATED; the same request truncates again at the same limit. (3) REFUSED; a refusal refuses again. (4) is the trap: the call is COMPLETED, because generation ended normally; the failure is in authorization or the file system, and the fix is a fresh observation, such as listing the directory. Step 3: the chapter's pitfall. A runtime that treats every failure like a dropped packet resends the identical request and repeats the cause. Expected errors: students mark row 4 as a failed call; the status describes generation, not the task. LINK: Lecture 12 owns the order of recovery decisions; Lecture 7 owns retrying an uncertain action safely.}
\end{frame}

% =============================================================================
% SECTION: CHECKS
% =============================================================================
\mlsysnote{[0.5 min] [CORE] Land the third segment: the status says whether a complete candidate arrived, and nothing else. Now: what does each check on a completed candidate actually prove?}
\mlsysdivider{What does each check on a candidate prove?}{The Foundation Model, \S2.6 and \S2.7}
\section{Checks}

\begin{frame}{Each check certifies one named property}
  \centering
  \mlsyssteps[width=0.98\textwidth,height=6.3cm,keepaspectratio]{ladder_step}{4}
  \mlsyscredit{The Foundation Model, table 3 and \S2.6.3.}
  \note{[2.5 min] [CORE] The chapter's four verification layers, ordered from cheap to costly. Step 1: a parser and schema validator certify fields and types; the read\_file call from our run passes, wrong path and all. Step 2: authorization checks the target against the task's grant at the moment of access, Saltzer and Schroeder's complete mediation; it says nothing about whether the read serves the goal. Step 3: a compiler or linter certifies what its analysis covers under its flags. Step 4: tests certify the inputs they exercised; a passing single-threaded test says nothing about the race in the connection pool. ENGAGE: an inspection call needs which rows before dispatch? (The first two; reading the file produces an observation, not a finished repair.) A patch needs all four, in a sandbox. WARN: ordering by cost saves compute, but each row tests a different property, not a tighter version of the same one. LINK: Lecture 1's closure evidence levels rank this same kind of evidence.}
\end{frame}

\begin{frame}{Self-review without new evidence lowered GPT-4's accuracy}
  \centering
  \mlsyssteps[width=0.90\textwidth,height=4.9cm,keepaspectratio]{huang_step}{2}\\[0.02cm]
  {\footnotesize Same model, asked to review and revise its own answer twice, with no new information.}
  \mlsyscredit{Huang et al., ``Large Language Models Cannot Self-Correct Reasoning Yet,'' ICLR 2024, Table 3 (intrinsic self-correction) and Table 2 (oracle labels), GPT-4 rows.}
  \note{[2 min] [OPTIONAL] Published measurements, read from the paper's tables. Step 1: GPT-4's accuracy with one answer (gray) and after two rounds of asking it to review and revise its own answer (red): 95.5 to 89.0 on GSM8K math, 82.0 to 80.0 on CommonSenseQA, 49.0 to 43.0 on HotpotQA. Self-review is one more generation from the same weights and the same context; a mistaken assumption usually survives it, and correct answers get changed to wrong ones. Step 2: green, the same loop stopped by the true label, an external oracle: 97.5, 85.5, and 59.0. External evidence is what helped, not the extra calls. ENGAGE: in an agent, what plays the role of the oracle? (A compiler, a test run, a file read; something outside the model.) WARN: the oracle numbers use labels a real agent does not have; they show where the gain comes from, not a method to deploy. LINK: learned judges are graded in Lecture 13.}
\end{frame}

\begin{frame}{How many passing patches are still defective?}
  \mlsysactivity{Estimate}{3 min}
  \begin{columns}[T]
    \begin{column}{0.36\textwidth}
      \begin{mlsyscard}{mlsysaccent}
        \small\textbf{Given (Sys.\ Persp.\ 2.1)}\\[0.05cm]
        Defective: 18\%\\
        Passes a defective patch: 0.65\\
        Passes a valid patch: 0.85
      \end{mlsyscard}
      {\footnotesize $\Pr(\text{Def}\mid\text{PASS}) = \dfrac{0.65 \times 0.18}{0.65 \times 0.18 + 0.85 \times 0.82}$}
    \end{column}
    \begin{column}{0.60\textwidth}
      \centering
      \mlsyssteps[width=\textwidth,height=3.6cm,keepaspectratio]{bayes_step}{3}\\[0.05cm]
      \uncover<3->{\footnotesize Rerunning the same check buys nothing: its errors repeat on the same patches.}
    \end{column}
  \end{columns}
  \mlsyscredit{The Foundation Model, Systems Perspective 2.1 and equation 4 (assumed rates). Arithmetic: images/calcs.py.}
  \note{[3 min] [CORE] ENGAGE: two minutes with a neighbor; of the patches that pass, what share is defective? Step 1: 18 percent of generated patches violate the concurrency property. Step 2: Bayes' rule (equation 4): 0.65 times 0.18 over that plus 0.85 times 0.82 gives 14.37 percent. The check moved the defect share only from 18 to 14.4 percent, because it rejects only 35 percent of defective patches while also rejecting 15 percent of valid ones. Step 3: the lever is the false-accept rate. Halving it to 0.325 lowers the residual to 7.7 percent. Repeating the same check does not, because its behavior on the survivors is correlated with the first pass. Expected errors: students answer 0.65 times 0.18 = 11.7 percent, the joint probability, forgetting to divide by all passes. WARN: every defective patch that passes also inflates the acceptance probability in $C_{\text{eff}}$. LINK: Lecture 3 owns verifier precision when choosing among many candidates.}
\end{frame}

\begin{frame}{A grammar mask keeps every legal token, including wrong paths}
  \centering
  \safeimg[width=0.96\textwidth,height=5.0cm,keepaspectratio]{mask.pdf}\\[0.05cm]
  {\small The mask sets illegal logits to $-\infty$, so a \textbf{completed} reply is valid syntax.}
  \mlsyscredit{Candidates and probabilities: the real call (images/run\_model.py). The Foundation Model, Definition 2.3, equation 5, and table 4.}
  \note{[2.5 min] [CORE] Grammar-constrained decoding (Definition 2.3): the engine tracks a parser state for the target grammar, and at each step adds a mask to the logits, 0 for tokens that keep a viable prefix and minus infinity for the rest (equation 5). Engines compile the schema once into an automaton so each step is a lookup. Applied to our real call at the first token of the path, the parser is inside a JSON string, so every one of the eight real candidates is legal, including "path", which leads to the invented file. Only raw control characters and bad escapes are removed. The mask moves the first row of the check ladder into generation and leaves the other rows where they were. ENGAGE: if the output limit cuts a masked reply at "src/po, may the runtime dispatch it? (No; it is TRUNCATED, and the mask never promotes an unfinished prefix.) WARN: a schema that requires a path and offers no "unknown" branch forces the model to invent one when the context lacks it; add an explicit unknown variant. FLEX: [CORE].}
\end{frame}

% =============================================================================
% SECTION: COST
% =============================================================================
\mlsysnote{[0.5 min] [CORE] Land the fourth segment: each check certifies one property, self-review adds none, and a grammar mask guarantees syntax only for a completed reply. Last: what does a call cost in time, and what does a trajectory of calls resend?}
\mlsysdivider{What sets a call's time, and what does a trajectory resend?}{The Foundation Model, \S2.8 and \S2.9}
\section{Cost}

\begin{frame}{Every decode step reads all the weights from memory}
  \begin{columns}[c]
    \begin{column}{0.42\textwidth}
      \centering\safeimg[width=\textwidth,height=4.4cm,keepaspectratio]{fiji-hbm-package.jpg}\\[0.02cm]
      {\footnotesize Accelerator die (center) and four stacks of high-bandwidth memory (HBM).}
    \end{column}
    \begin{column}{0.55\textwidth}
      \small
      $\tau_{\text{prefill}} \approx \dfrac{2P}{R_{\text{peak}}} = \mathbf{71.35}\ \mu\text{s}$ per prompt token\\[0.2cm]
      $\tau_{\text{decode}} \ge \dfrac{P\,b}{BW} = \mathbf{21.075}$ ms per output token\\[0.25cm]
      \uncover<2->{{\large\bfseries\textcolor{mlsysaccent}{One output token $\approx$ 295 prompt tokens}}\\[0.1cm]
      {\footnotesize Reference: 70.6B parameters $P$, 8-bit ($b$ = 1 byte), $R_{\text{peak}}$ = 1,979 TFLOP/s, $BW$ = 3.35 TB/s.}}
    \end{column}
  \end{columns}
  \mlsyscredit{Photo: AMD Fiji package, C. Spille/pcgameshardware.de, CC BY-SA 4.0, Wikimedia Commons. Floors: The Foundation Model, equations 7 and 8, \S2.8.2.}
  \note{[3 min] [CORE] The photo shows an accelerator package: the large compute die in the center and four stacks of high-bandwidth memory beside it. The weights live in those stacks. Prefill reads the whole prompt in one parallel pass, so each weight is fetched once and reused across thousands of prompt tokens: about $2P$ operations per token at the peak rate, 71.35 microseconds per token for the chapter's reference deployment. Decode produces one token per pass at batch one, so every output token reads every weight: 70.6 GB over 3.35 TB/s, at least 21.075 ms. Step 2: the ratio (equation 8) is about 295, independent of model size. This is output-bound latency, the Part I principle. The chapter's first repair turn makes it concrete: 8,000 prompt tokens take 0.571 s, 200 output tokens take 4.215 s; output is 2.4 percent of the tokens and 88.1 percent of the time. WARN: these are floors for unbatched decode, not a measured service; batching amortizes the weight read (the Inference Cost Foundations appendix derives the batched bound).}
\end{frame}

\begin{frame}{Real traffic is mostly input tokens; time goes to output}
  \begin{columns}[T]
    \begin{column}{0.48\textwidth}
      \centering{\small\textbf{Measured on a laptop: our 0.5B model}}\\[0.05cm]
      \safeimg[width=\textwidth,height=3.3cm,keepaspectratio]{laptop.pdf}\\[0.05cm]
      {\footnotesize 2,048-token prefill in 2.05 s; decode 35.1 ms per token. One output token $\approx$ 35 prompt tokens.}
    \end{column}
    \begin{column}{0.48\textwidth}
      \centering{\small\textbf{Code-completion trace, 8,819 requests}}\\[0.05cm]
      \safeimg[width=\textwidth,height=3.3cm,keepaspectratio]{azure.pdf}\\[0.05cm]
      {\footnotesize Input is 98.7\% of all tokens. At the reference floors, output would take 80\% of the time (derived).}
    \end{column}
  \end{columns}
  \mlsyscredit{Left: images/run\_model.py, CPU float32, batch 1, medians of 5 and 3 runs. Right: Azure LLM inference trace 2023, code (Patel et al., ISCA 2024), CC BY 4.0.}
  \note{[2 min] [OPTIONAL] Two real measurements. Left: our 0.5B model on a laptop CPU, batch one. A 2,048-token prompt prefilled in 2.05 s, about 1.0 ms per token; decode took 35.1 ms per token. One output token cost about 35 prompt tokens of time (25 to 41 across runs). The ratio is smaller than 295 because a laptop CPU does not reach its arithmetic peak on prefill; the direction is the same. Right: a public trace of 8,819 production code-completion requests to an LLM service. The median request sent 1,469 prompt tokens and received 13 output tokens; input is 98.7 percent of all tokens. If those requests ran at the chapter's reference floors, output would still take 80 percent of the service time (derived). ENGAGE: an engineer wants to cut latency by 10 percent; trim the prompt or the reply? (The reply, unless queueing or a very long prompt dominates.) WARN: the trace is code completion, not an agent; it shows the token mix of real traffic.}
\end{frame}

\begin{frame}{How much input does a 30-turn repair resend?}
  \mlsysactivity{Predict}{3.5 min}
  \begin{columns}[c]
    \begin{column}{0.36\textwidth}
      \begin{mlsyscard}{mlsysaccent}
        \small\textbf{Given (Napkin Math 2.2)}\\[0.05cm]
        First prompt: 8,000\\
        Appended per turn: 1,500\\
        Output per call: 200\\
        Turns $H$: 30
      \end{mlsyscard}
      \vspace{0.1cm}
      {\small\choice{A}{about 50k}\\ \choice{B}{about 250k}\\ \answerat{2}{C}{about 900k}}
    \end{column}
    \begin{column}{0.60\textwidth}
      \centering
      \mlsyssteps[width=\textwidth,height=4.2cm,keepaspectratio]{resend_step}{3}\\[0.02cm]
      \uncover<3->{\footnotesize Quadratic resend: 892,500 tokens. Prefix reuse: 51,500 fresh, \textbf{17.3$\times$} less.}
    \end{column}
  \end{columns}
  \mlsyscredit{The Foundation Model, Napkin Math 2.2 and equation 9 (scenario values). Arithmetic: images/calcs.py.}
  \note{[3.5 min] [CORE] ENGAGE: the runtime keeps an append-only log and resends the full history on every call. Vote on the total input tokens sent over 30 turns. Step 2: C. Equation 9: 30 times 8,000 plus 1,500 times 30 times 29 over 2 is 892,500 tokens. Input is 99.3 percent of all tokens; output is only 6,000. Students who pick A add the appends once; students who pick B forget that every turn resends everything before it. The growth is quadratic in the horizon. Step 3: with prefix reuse, the server keeps the attention state of a repeated prefix, so only the new suffix needs fresh prefill: 51,500 tokens, 17.3 times less. At the reference floor prefill work falls from 63.68 s to 3.67 s; decode stays 126.45 s. Doubling to 60 turns multiplies resent input 3.5 times but fresh prefill only 1.9 times. WARN: reuse requires an exact token prefix; changing one early token, such as a timestamp in the system prompt, forces fresh prefill from there. LINK: Lecture 4 orders context for reuse; Lecture 5 owns the cache; Lecture 19 prices it.}
\end{frame}

\begin{frame}{A whole-file edit misses a 15 s deadline that a diff meets}
  \begin{columns}[c]
    \begin{column}{0.62\textwidth}
      \centering
      \mlsyssteps[width=\textwidth,height=4.6cm,keepaspectratio]{deadline_step}{2}
    \end{column}
    \begin{column}{0.35\textwidth}
      \small Whole file: 1,850 output tokens.\\ Local edit: 80.\\[0.15cm]
      \uncover<2->{The edit is 17.5$\times$ faster, and 57\% of trials end accepted ($0.95 \times 0.60$).\\[0.1cm]
      {\footnotesize Every trial stays in the denominator.}}
    \end{column}
  \end{columns}
  \mlsyscredit{The Foundation Model, Napkin Math 2.3 and table 5 (scenario values, reference floors). Arithmetic: images/calcs.py.}
  \note{[3 min] [CORE] Two output interfaces for the same repair. Step 1: a whole-file replacement adds 350 instruction tokens to the 8,000-token context and generates 1,850 output tokens: 39.58 s at the reference floors, 98.5 percent of it decode, 24.6 s past the 15 s deadline. Step 2: a localized edit adds 120 instruction tokens and generates 80: 2.27 s, 17.5 times faster, with a 6.6 times margin. But a local edit must find its anchor text: in the chapter's hypothetical, 95 percent apply and 60 percent of those pass the tests, so 57 percent of all trials yield an accepted repair. That 57 percent is the acceptance probability in $C_{\text{eff}}$. ENGAGE: if the deadline were 60 s, which format wins? (Then applicability decides; a whole file that always applies can beat a diff that misses stale anchors.) WARN: a fair comparison fixes the model, tools, tests, limits, and retry policy, and keeps every timed-out or unparsable trial in the denominator. LINK: Lecture 13 owns matched-budget evaluation of whole agents.}
\end{frame}

% =============================================================================
% CLOSE
% =============================================================================
\begin{frame}{Which records would have caught the bugs?}
  \mlsysactivity{Retrieval}{3 min}
  \begin{columns}[c]
    \begin{column}{0.36\textwidth}
      \centering\safeimg[width=\textwidth,height=2.6cm,keepaspectratio]{postmortem_timeline.png}\\[0.05cm]
      {\footnotesize Requests unchanged; outputs corrupted for weeks.}
    \end{column}
    \begin{column}{0.61\textwidth}
      \only<1|handout:0>{\normalsize Close your notes. For each runtime record, say what it would have shown during the 2025 incident.}%
      \only<2->{\footnotesize
      \begin{tabular}{@{}P{3.0cm}P{5.0cm}@{}}\toprule
        \textbf{Record} & \textbf{What it shows} \\\midrule
        Status envelope & nothing (replies were COMPLETED) \\
        Schema check & nothing (replies parsed) \\
        Serving config log & which path served each call \\
        External tests & the wrong tokens, in code \\
        Recorded responses & replay a failure exactly \\\bottomrule
      \end{tabular}}
    \end{column}
  \end{columns}
  \mlsyscredit{Timeline: Anthropic 2025 postmortem. The Foundation Model, War Story 2.1 and \S2.4.3.}
  \note{[3 min] [CORE] Back to the hook. Step 1: retrieval, 60 seconds. Step 2: the status envelope and the schema check, the two cheapest records, show nothing, because the corrupted replies finished normally and parsed. That is the fail-plausible fault. A log of model revision and serving configuration per call shows which path served each failing call, which is how the routing bug could be traced. External tests, the compiler and test runner in the repair loop, see the wrong token in the code. Recorded responses let an engineer replay the failure exactly, since resending the request lands on a different path. The provider's own fix added continuous canary evaluations on production traffic, an external check at the provider's side. ENGAGE: of the answers you gave at the start, which were right? WARN: none of these makes the model reliable; they make its failures visible to the runtime, which is the point of the contract.}
\end{frame}

\begin{frame}{Takeaways}
  \thumbline{thumb_boundary_step3.pdf}{A model call returns a proposal; the runtime validates, authorizes, executes, and observes.}\\[0.12cm]
  \thumbline{thumb_probs_step3.pdf}{Sampling reshapes candidates without testing them; only a recorded response replays exactly.}\\[0.12cm]
  \thumbline{thumb_rmrf.pdf}{Read the status before the payload: only COMPLETED enters the candidate path.}\\[0.12cm]
  \thumbline{thumb_ladder_step4.pdf}{Each check certifies one property; self-review adds none, and a mask guarantees only syntax.}\\[0.12cm]
  \thumbline{thumb_resend_step3.pdf}{One output token costs about 295 prompt tokens; resent history grows with the square of the horizon.}
  \note{[1 min] [CORE] One line per segment. Ask students which slide taught each line. If short on time, read the first and the last.}
\end{frame}

\begin{frame}{Open questions}
  \normalsize
  \begin{itemize}\setlength{\itemsep}{0.3cm}
    \item A provider's serving path can change under a pinned model name. What evidence should a provider publish so an agent runtime can tell?
    \item Grammar masks remove syntax errors at a per-token cost. When should a runtime prefer streaming rejection instead?
    \item Reasoning models spend hidden output tokens. Should the runtime or the provider set that budget?
  \end{itemize}
  \note{[1 min] [OPTIONAL] No single answer. The first follows from War Story 2.1 and \S2.5.1 on mutable aliases and pinned revisions. The second is the chapter's Exercise 6 as a discussion: compare added latency per proposal under post-hoc validation, streaming rejection, and masking. The third follows from \S2.3.2: hidden reasoning tokens debit the output reserve, so they change both admission and cost. Assign one as a short written response.}
\end{frame}

\begin{frame}{Read and explore}
  \small
  \begin{itemize}\setlength{\itemsep}{0.02cm}
    \item \textbf{Read:} The Foundation Model, \S2.2 to \S2.9; Checkpoints 2.1 to 2.4; Napkin Math 2.1 to 2.3 by hand; Exercises 1, 4, and 9.
    \item Birrell and Nelson, ``Implementing Remote Procedure Calls,'' 1984: the unknown outcome.
    \item Holtzman et al., ``The Curious Case of Neural Text Degeneration,'' ICLR 2020: nucleus sampling.
    \item Willard and Louf, ``Efficient Guided Generation for LLMs,'' 2023: grammar masks as automata.
    \item Pope et al., ``Efficiently Scaling Transformer Inference,'' MLSys 2023: prefill and decode floors.
    \item Huang et al., ``LLMs Cannot Self-Correct Reasoning Yet,'' ICLR 2024.
  \end{itemize}
  \note{[0.5 min] [CORE] The chapter first, with its checkpoints and the three napkin-math notebooks. Exercise 9 transfers the truncation rule to the remediation agent, where a fragment like --replicas=1 changes production. Birrell and Nelson explain why a timeout leaves the outcome unknown. Holtzman et al. introduce top-p. Willard and Louf compile grammars to token masks. Pope et al. derive the two serving floors. Huang et al. measure self-correction. FLEX: run images/run\_model.py in a lab to let students see their own model's wire text and logits.}
\end{frame}

\begin{frame}[c]
  \begin{columns}[c]
    \begin{column}{0.34\textwidth}
      \centering\mlsyskitfig[width=\textwidth,height=4.2cm,keepaspectratio]{stack_model}
    \end{column}
    \begin{column}{0.62\textwidth}
      {\small\textcolor{lightgray}{NEXT LECTURE \textbar{} TEST-TIME COMPUTE}}\\[0.5cm]
      {\LARGE\bfseries When is another token, candidate, check, or model call worth its cost?}
    \end{column}
  \end{columns}
  \note{[0.5 min] [CORE] Today ended at one call and one candidate. A completed candidate can still rest on the wrong diagnosis. Lecture 3 asks how to spend more before committing, through a longer derivation, more sampled candidates, or a targeted test, and why only the test adds a fact the model did not have.}
\end{frame}

\end{document}
